%=======================================================================
%  FIGURE -- how a coefficient drift score is computed
%
%  Include from a chapter with:
%      \begin{figure}[htbp]
%        \centering
%        \resizebox{0.82\textwidth}{!}{%=======================================================================
%  FIGURE -- how a coefficient drift score is computed
%
%  Include from a chapter with:
%      \begin{figure}[htbp]
%        \centering
%        \resizebox{0.82\textwidth}{!}{%=======================================================================
%  FIGURE -- how a coefficient drift score is computed
%
%  Include from a chapter with:
%      \begin{figure}[htbp]
%        \centering
%        \resizebox{0.82\textwidth}{!}{%=======================================================================
%  FIGURE -- how a coefficient drift score is computed
%
%  Include from a chapter with:
%      \begin{figure}[htbp]
%        \centering
%        \resizebox{0.82\textwidth}{!}{\input{figures/drift_pipeline}}
%        \caption{...}
%        \label{fig:drift-pipeline}
%      \end{figure}
%
%  Why this figure exists.
%
%  Equations 3.1 to 3.4 give the recursive least squares update: how one
%  coefficient moves from one timestep to the next. They do not show the
%  thing the detector is built on, which is that two passes are run over
%  two disjoint windows and their coefficients are differenced. A reader
%  can follow every equation in Chapter 3 and still not know that w^(0)
%  comes from an offline fit on the training split while w(t) comes from
%  an online pass over the evaluation split, nor that the ratio is taken
%  per coefficient before the maximum is taken across them.
%
%  STYLE NAMES, because the first version of this figure would not
%  compile. Every style below is prefixed `dp'. The first version named
%  one of them `out', which is a reserved TikZ key -- it is the outgoing
%  angle in `to[out=30,in=150]' -- so `\node[out, ...]' was read as the
%  key /tikz/out with no value, and the rest of the picture unravelled
%  from there. `in', `at', `to', `pos', `scale' and `shift' are reserved
%  in the same way. A prefix costs nothing and rules the whole class of
%  collisions out.
%
%  LAYOUT: placement is relative throughout, via positioning's below=of
%  and right=of, so the boxes cannot overlap however the text inside them
%  reflows.
%=======================================================================
\begin{tikzpicture}[
  font=\small,
  dpbox/.style={draw=black!45, rounded corners=2pt, align=center,
                text width=54mm, inner sep=2.2mm, font=\footnotesize},
  dpsrc/.style={dpbox, fill=black!4},
  dpmid/.style={dpbox, text width=74mm},
  dpend/.style={dpbox, text width=74mm, draw=black!65, very thick},
  dparr/.style={-{Stealth[length=2mm]}, draw=black!60},
  dplbl/.style={font=\scriptsize\itshape, text=black!60, inner sep=1pt},
  node distance=7mm
]

% ---------------------------------------------- the two windows
\node[dpsrc] (train) {\textbf{fault-free training window}\\ first 70\%
  of the recording, no fault present};
\node[dpsrc, right=14mm of train] (evalwin) {\textbf{evaluation window}\\ last 30
 \%, containing the injected faults};

% ---------------------------------------------- the two passes
\node[dpbox, below=of train] (offline)
  {\textbf{offline fit}\\ least squares over the whole window gives one baseline
   coefficient vector per target\\[1mm] $\mathbf{w}^{(0)}_{v}$};
\node[dpbox, below=of evalwin] (online)
  {\textbf{online pass}\\ the RLS filter of Equations~\ref{eq:rls-gain}
   to~\ref{eq:rls-cov} updates the coefficients at every
   timestep\\[1mm] $\mathbf{w}_{v}(t)$\\[0.8mm]\scriptsize\itshape first 15 steps discarded as warm-up};

\draw[dparr] (train) -- node[dplbl, right] {fit once} (offline);
\draw[dparr] (evalwin) -- node[dplbl, right] {track continuously} (online);

% ---------------------------------------------- the difference
\coordinate (dpmerge) at ($(offline.south)!0.5!(online.south)$);
\node[dpmid, below=11mm of dpmerge] (diff)
  {\textbf{departure of one coefficient}\\[0.6mm]
   $d_{v,k}(t) \;=\; w_{v,k}(t) \;-\; w^{(0)}_{v,k}$\\[0.6mm]
   the modelled relationship has moved, not the signal};

\draw[dparr] (offline) -- (diff);
\draw[dparr] (online) -- (diff);


% ---------------------------------------------- normalisation
\node[dpmid, below=of diff] (norm)
  {\textbf{divided by that same coefficient's own fault-free drift}\\[0.6mm]
   $\bigl|d_{v,k}(t)\bigr| \;\big/\;
    \kappa\,\bigl\lVert d^{\mathrm{base}}_{v,k}\bigr\rVert_2$\\[0.6mm]
   which is what lets coefficients of different scale be compared};
\draw[dparr] (diff) -- (norm);

% ---------------------------------------------- the maximum
\node[dpmid, below=of norm] (agg)
  {\textbf{maximum over every coefficient of every target}\\[0.6mm]
   $s(t) \;=\; \displaystyle\max_{v,\,k}$ of the ratio above\\[0.6mm]
   one broken relationship is enough to raise the score};
\draw[dparr] (norm) -- (agg);

% ---------------------------------------------- the decision
\node[dpend, below=of agg] (flag)
  {\textbf{decision}\\[0.6mm]
   the continuous $s(t)$ is what AUC-ROC is computed over; the reported
   operating point flags a sample when $s(t)$ exceeds the 95th percentile
   of the fault-free score distribution, fixed before any label is seen};
\draw[dparr] (agg) -- (flag);

\end{tikzpicture}
}
%        \caption{...}
%        \label{fig:drift-pipeline}
%      \end{figure}
%
%  Why this figure exists.
%
%  Equations 3.1 to 3.4 give the recursive least squares update: how one
%  coefficient moves from one timestep to the next. They do not show the
%  thing the detector is built on, which is that two passes are run over
%  two disjoint windows and their coefficients are differenced. A reader
%  can follow every equation in Chapter 3 and still not know that w^(0)
%  comes from an offline fit on the training split while w(t) comes from
%  an online pass over the evaluation split, nor that the ratio is taken
%  per coefficient before the maximum is taken across them.
%
%  STYLE NAMES, because the first version of this figure would not
%  compile. Every style below is prefixed `dp'. The first version named
%  one of them `out', which is a reserved TikZ key -- it is the outgoing
%  angle in `to[out=30,in=150]' -- so `\node[out, ...]' was read as the
%  key /tikz/out with no value, and the rest of the picture unravelled
%  from there. `in', `at', `to', `pos', `scale' and `shift' are reserved
%  in the same way. A prefix costs nothing and rules the whole class of
%  collisions out.
%
%  LAYOUT: placement is relative throughout, via positioning's below=of
%  and right=of, so the boxes cannot overlap however the text inside them
%  reflows.
%=======================================================================
\begin{tikzpicture}[
  font=\small,
  dpbox/.style={draw=black!45, rounded corners=2pt, align=center,
                text width=54mm, inner sep=2.2mm, font=\footnotesize},
  dpsrc/.style={dpbox, fill=black!4},
  dpmid/.style={dpbox, text width=74mm},
  dpend/.style={dpbox, text width=74mm, draw=black!65, very thick},
  dparr/.style={-{Stealth[length=2mm]}, draw=black!60},
  dplbl/.style={font=\scriptsize\itshape, text=black!60, inner sep=1pt},
  node distance=7mm
]

% ---------------------------------------------- the two windows
\node[dpsrc] (train) {\textbf{fault-free training window}\\ first 70\%
  of the recording, no fault present};
\node[dpsrc, right=14mm of train] (evalwin) {\textbf{evaluation window}\\ last 30
 \%, containing the injected faults};

% ---------------------------------------------- the two passes
\node[dpbox, below=of train] (offline)
  {\textbf{offline fit}\\ least squares over the whole window gives one baseline
   coefficient vector per target\\[1mm] $\mathbf{w}^{(0)}_{v}$};
\node[dpbox, below=of evalwin] (online)
  {\textbf{online pass}\\ the RLS filter of Equations~\ref{eq:rls-gain}
   to~\ref{eq:rls-cov} updates the coefficients at every
   timestep\\[1mm] $\mathbf{w}_{v}(t)$\\[0.8mm]\scriptsize\itshape first 15 steps discarded as warm-up};

\draw[dparr] (train) -- node[dplbl, right] {fit once} (offline);
\draw[dparr] (evalwin) -- node[dplbl, right] {track continuously} (online);

% ---------------------------------------------- the difference
\coordinate (dpmerge) at ($(offline.south)!0.5!(online.south)$);
\node[dpmid, below=11mm of dpmerge] (diff)
  {\textbf{departure of one coefficient}\\[0.6mm]
   $d_{v,k}(t) \;=\; w_{v,k}(t) \;-\; w^{(0)}_{v,k}$\\[0.6mm]
   the modelled relationship has moved, not the signal};

\draw[dparr] (offline) -- (diff);
\draw[dparr] (online) -- (diff);


% ---------------------------------------------- normalisation
\node[dpmid, below=of diff] (norm)
  {\textbf{divided by that same coefficient's own fault-free drift}\\[0.6mm]
   $\bigl|d_{v,k}(t)\bigr| \;\big/\;
    \kappa\,\bigl\lVert d^{\mathrm{base}}_{v,k}\bigr\rVert_2$\\[0.6mm]
   which is what lets coefficients of different scale be compared};
\draw[dparr] (diff) -- (norm);

% ---------------------------------------------- the maximum
\node[dpmid, below=of norm] (agg)
  {\textbf{maximum over every coefficient of every target}\\[0.6mm]
   $s(t) \;=\; \displaystyle\max_{v,\,k}$ of the ratio above\\[0.6mm]
   one broken relationship is enough to raise the score};
\draw[dparr] (norm) -- (agg);

% ---------------------------------------------- the decision
\node[dpend, below=of agg] (flag)
  {\textbf{decision}\\[0.6mm]
   the continuous $s(t)$ is what AUC-ROC is computed over; the reported
   operating point flags a sample when $s(t)$ exceeds the 95th percentile
   of the fault-free score distribution, fixed before any label is seen};
\draw[dparr] (agg) -- (flag);

\end{tikzpicture}
}
%        \caption{...}
%        \label{fig:drift-pipeline}
%      \end{figure}
%
%  Why this figure exists.
%
%  Equations 3.1 to 3.4 give the recursive least squares update: how one
%  coefficient moves from one timestep to the next. They do not show the
%  thing the detector is built on, which is that two passes are run over
%  two disjoint windows and their coefficients are differenced. A reader
%  can follow every equation in Chapter 3 and still not know that w^(0)
%  comes from an offline fit on the training split while w(t) comes from
%  an online pass over the evaluation split, nor that the ratio is taken
%  per coefficient before the maximum is taken across them.
%
%  STYLE NAMES, because the first version of this figure would not
%  compile. Every style below is prefixed `dp'. The first version named
%  one of them `out', which is a reserved TikZ key -- it is the outgoing
%  angle in `to[out=30,in=150]' -- so `\node[out, ...]' was read as the
%  key /tikz/out with no value, and the rest of the picture unravelled
%  from there. `in', `at', `to', `pos', `scale' and `shift' are reserved
%  in the same way. A prefix costs nothing and rules the whole class of
%  collisions out.
%
%  LAYOUT: placement is relative throughout, via positioning's below=of
%  and right=of, so the boxes cannot overlap however the text inside them
%  reflows.
%=======================================================================
\begin{tikzpicture}[
  font=\small,
  dpbox/.style={draw=black!45, rounded corners=2pt, align=center,
                text width=54mm, inner sep=2.2mm, font=\footnotesize},
  dpsrc/.style={dpbox, fill=black!4},
  dpmid/.style={dpbox, text width=74mm},
  dpend/.style={dpbox, text width=74mm, draw=black!65, very thick},
  dparr/.style={-{Stealth[length=2mm]}, draw=black!60},
  dplbl/.style={font=\scriptsize\itshape, text=black!60, inner sep=1pt},
  node distance=7mm
]

% ---------------------------------------------- the two windows
\node[dpsrc] (train) {\textbf{fault-free training window}\\ first 70\%
  of the recording, no fault present};
\node[dpsrc, right=14mm of train] (evalwin) {\textbf{evaluation window}\\ last 30
 \%, containing the injected faults};

% ---------------------------------------------- the two passes
\node[dpbox, below=of train] (offline)
  {\textbf{offline fit}\\ least squares over the whole window gives one baseline
   coefficient vector per target\\[1mm] $\mathbf{w}^{(0)}_{v}$};
\node[dpbox, below=of evalwin] (online)
  {\textbf{online pass}\\ the RLS filter of Equations~\ref{eq:rls-gain}
   to~\ref{eq:rls-cov} updates the coefficients at every
   timestep\\[1mm] $\mathbf{w}_{v}(t)$\\[0.8mm]\scriptsize\itshape first 15 steps discarded as warm-up};

\draw[dparr] (train) -- node[dplbl, right] {fit once} (offline);
\draw[dparr] (evalwin) -- node[dplbl, right] {track continuously} (online);

% ---------------------------------------------- the difference
\coordinate (dpmerge) at ($(offline.south)!0.5!(online.south)$);
\node[dpmid, below=11mm of dpmerge] (diff)
  {\textbf{departure of one coefficient}\\[0.6mm]
   $d_{v,k}(t) \;=\; w_{v,k}(t) \;-\; w^{(0)}_{v,k}$\\[0.6mm]
   the modelled relationship has moved, not the signal};

\draw[dparr] (offline) -- (diff);
\draw[dparr] (online) -- (diff);


% ---------------------------------------------- normalisation
\node[dpmid, below=of diff] (norm)
  {\textbf{divided by that same coefficient's own fault-free drift}\\[0.6mm]
   $\bigl|d_{v,k}(t)\bigr| \;\big/\;
    \kappa\,\bigl\lVert d^{\mathrm{base}}_{v,k}\bigr\rVert_2$\\[0.6mm]
   which is what lets coefficients of different scale be compared};
\draw[dparr] (diff) -- (norm);

% ---------------------------------------------- the maximum
\node[dpmid, below=of norm] (agg)
  {\textbf{maximum over every coefficient of every target}\\[0.6mm]
   $s(t) \;=\; \displaystyle\max_{v,\,k}$ of the ratio above\\[0.6mm]
   one broken relationship is enough to raise the score};
\draw[dparr] (norm) -- (agg);

% ---------------------------------------------- the decision
\node[dpend, below=of agg] (flag)
  {\textbf{decision}\\[0.6mm]
   the continuous $s(t)$ is what AUC-ROC is computed over; the reported
   operating point flags a sample when $s(t)$ exceeds the 95th percentile
   of the fault-free score distribution, fixed before any label is seen};
\draw[dparr] (agg) -- (flag);

\end{tikzpicture}
}
%        \caption{...}
%        \label{fig:drift-pipeline}
%      \end{figure}
%
%  Why this figure exists.
%
%  Equations 3.1 to 3.4 give the recursive least squares update: how one
%  coefficient moves from one timestep to the next. They do not show the
%  thing the detector is built on, which is that two passes are run over
%  two disjoint windows and their coefficients are differenced. A reader
%  can follow every equation in Chapter 3 and still not know that w^(0)
%  comes from an offline fit on the training split while w(t) comes from
%  an online pass over the evaluation split, nor that the ratio is taken
%  per coefficient before the maximum is taken across them.
%
%  STYLE NAMES, because the first version of this figure would not
%  compile. Every style below is prefixed `dp'. The first version named
%  one of them `out', which is a reserved TikZ key -- it is the outgoing
%  angle in `to[out=30,in=150]' -- so `\node[out, ...]' was read as the
%  key /tikz/out with no value, and the rest of the picture unravelled
%  from there. `in', `at', `to', `pos', `scale' and `shift' are reserved
%  in the same way. A prefix costs nothing and rules the whole class of
%  collisions out.
%
%  LAYOUT: placement is relative throughout, via positioning's below=of
%  and right=of, so the boxes cannot overlap however the text inside them
%  reflows.
%=======================================================================
\begin{tikzpicture}[
  font=\small,
  dpbox/.style={draw=black!45, rounded corners=2pt, align=center,
                text width=54mm, inner sep=2.2mm, font=\footnotesize},
  dpsrc/.style={dpbox, fill=black!4},
  dpmid/.style={dpbox, text width=74mm},
  dpend/.style={dpbox, text width=74mm, draw=black!65, very thick},
  dparr/.style={-{Stealth[length=2mm]}, draw=black!60},
  dplbl/.style={font=\scriptsize\itshape, text=black!60, inner sep=1pt},
  node distance=7mm
]

% ---------------------------------------------- the two windows
\node[dpsrc] (train) {\textbf{fault-free training window}\\ first 70\%
  of the recording, no fault present};
\node[dpsrc, right=14mm of train] (evalwin) {\textbf{evaluation window}\\ last 30
 \%, containing the injected faults};

% ---------------------------------------------- the two passes
\node[dpbox, below=of train] (offline)
  {\textbf{offline fit}\\ least squares over the whole window gives one baseline
   coefficient vector per target\\[1mm] $\mathbf{w}^{(0)}_{v}$};
\node[dpbox, below=of evalwin] (online)
  {\textbf{online pass}\\ the RLS filter of Equations~\ref{eq:rls-gain}
   to~\ref{eq:rls-cov} updates the coefficients at every
   timestep\\[1mm] $\mathbf{w}_{v}(t)$\\[0.8mm]\scriptsize\itshape first 15 steps discarded as warm-up};

\draw[dparr] (train) -- node[dplbl, right] {fit once} (offline);
\draw[dparr] (evalwin) -- node[dplbl, right] {track continuously} (online);

% ---------------------------------------------- the difference
\coordinate (dpmerge) at ($(offline.south)!0.5!(online.south)$);
\node[dpmid, below=11mm of dpmerge] (diff)
  {\textbf{departure of one coefficient}\\[0.6mm]
   $d_{v,k}(t) \;=\; w_{v,k}(t) \;-\; w^{(0)}_{v,k}$\\[0.6mm]
   the modelled relationship has moved, not the signal};

\draw[dparr] (offline) -- (diff);
\draw[dparr] (online) -- (diff);


% ---------------------------------------------- normalisation
\node[dpmid, below=of diff] (norm)
  {\textbf{divided by that same coefficient's own fault-free drift}\\[0.6mm]
   $\bigl|d_{v,k}(t)\bigr| \;\big/\;
    \kappa\,\bigl\lVert d^{\mathrm{base}}_{v,k}\bigr\rVert_2$\\[0.6mm]
   which is what lets coefficients of different scale be compared};
\draw[dparr] (diff) -- (norm);

% ---------------------------------------------- the maximum
\node[dpmid, below=of norm] (agg)
  {\textbf{maximum over every coefficient of every target}\\[0.6mm]
   $s(t) \;=\; \displaystyle\max_{v,\,k}$ of the ratio above\\[0.6mm]
   one broken relationship is enough to raise the score};
\draw[dparr] (norm) -- (agg);

% ---------------------------------------------- the decision
\node[dpend, below=of agg] (flag)
  {\textbf{decision}\\[0.6mm]
   the continuous $s(t)$ is what AUC-ROC is computed over; the reported
   operating point flags a sample when $s(t)$ exceeds the 95th percentile
   of the fault-free score distribution, fixed before any label is seen};
\draw[dparr] (agg) -- (flag);

\end{tikzpicture}
